\section{Task 2: Classification and Hard-Routed Specialists}\label{sec:t2}

\begin{figure*}[t]
\centering
\includegraphics[width=0.6\textwidth]{arch_task2_task3.pdf}
\caption{Top: hard routing (Task~2); the classifier output $r$ selects a specialist or the identity bypass. Bottom: the soft mixture of experts (Task~3) built from the same components.}\label{fig:arch23}
\end{figure*}

\textbf{Classifier.} Four stride-2 stages of two $3\times3$ convolutions (16--128 channels, batch norm, ReLU, Dropout2d), global average pooling, dropout and a linear layer give $p=[p_{clean},p_{salt},p_{blur},p_{occ}]$ with 0.295\,M parameters and $r=\arg\max_k p_k$. It is trained with multiclass cross-entropy on exactly balanced batches; Optuna tuned learning rate, batch size, channel configuration, dropout and weight decay (objective: validation macro-F1), selecting learning rate $1.5\times10^{-3}$, batch 32, dropout 0.15, weight decay $1.3\times10^{-4}$, 30 epochs. On the test protocol it reaches accuracy \nClfAcc{} and macro-F1 \nClfMacroF{} (Table~\ref{tab:classifier}, Fig.~\ref{fig:cm}). Salt-and-pepper is recognised perfectly; the \nMisrouteCount{} errors (\nMisroute\%) are clean images taken for blurred or occluded ones, and low-severity occlusions taken for clean images.

\begin{table}[t]
\centering
\caption{Classifier metrics on the test protocol.}\label{tab:classifier}
\small
\adjustbox{max width=\columnwidth}{\begin{tabular}{lrccc}
\toprule
Class & Support & Precision & Recall & F1 \\
\midrule
clean & 3,669 & 0.965 & 0.977 & 0.971 \\
salt-and-pepper & 11,007 & 1.000 & 1.000 & 1.000 \\
Gaussian blur & 11,007 & 0.994 & 1.000 & 0.997 \\
occlusion & 11,007 & 0.998 & 0.988 & 0.993 \\
\midrule
macro average & 36,690 & 0.989 & 0.991 & 0.990 \\
\multicolumn{2}{l}{overall accuracy} & \multicolumn{3}{c}{0.9941} \\
\bottomrule
\end{tabular}
}
\end{table}

\begin{figure}[t]
\centering
\includegraphics[width=0.75\columnwidth]{confusion_matrix.pdf}
\caption{Row-normalised confusion matrix of the classifier.}\label{fig:cm}
\end{figure}

\textbf{Specialists.} Three autoencoders $A_{salt}$, $A_{blur}$, $A_{occ}$ share the Task~1 architecture but have independent parameters, each trained only on its own corruption with the same clean target. As the assignment allows, one shared Optuna search (each trial trains all three for six epochs on 1,000 images) selects a common configuration (learning rate $6.8\times10^{-4}$, batch 32, dropout 0.07, $\alpha=0.60$); the specialists are then trained for 100 epochs. Inference is
\begin{equation}
\hat{x}=\begin{cases}\tilde{x}, & r=\text{clean},\\ A_{salt}(\tilde{x}), & r=\text{salt},\\ A_{blur}(\tilde{x}), & r=\text{blur},\\ A_{occ}(\tilde{x}), & r=\text{occlusion},\end{cases}
\end{equation}
so clean inputs take an identity bypass. \emph{Oracle} routing (true label) measures the specialists; \emph{predicted} routing measures the deployable system.

\textbf{Colour loss and its correction.} The first specialists (Optuna chose $\alpha=0.29$) produced washed-out, sepia-toned images, which we noticed when inspecting the outputs. Colourfulness measured on \nColourN{} salt-and-pepper test inputs confirmed it (Table~\ref{tab:colour}): the first specialist kept \nColSpSpecFirst\% of the colour and the first mixture \nColSpMoeFirst\%, the universal model \nColSpUniversal\%. SSIM is weakly sensitive to a global chroma shift, so neither the loss nor the Optuna objective penalised it. We restricted $\alpha$ to $[0.55,0.95]$ and retrained the specialists and the mixture: the final specialist keeps \nColSpSpecFinal\% and the mixture \nColSpMoeFinal\%; the L1 error of the routed system fell by \nLOneDropPct\% (\nHardPredLOneFirst{} to \nHardPredLOneFinal{}) with SSIM essentially unchanged. The occlusion specialist still keeps slightly less colour than the universal model (\nColOcSpecFinal\% against \nColOcUniversal\%).

\begin{table}[t]
\centering
\caption{Colour preserved (colourfulness of the output relative to the clean image) on \nColourN{} salt-and-pepper and occlusion test inputs.}\label{tab:colour}
\small
\adjustbox{max width=\columnwidth}{\begin{tabular}{lcc}
\toprule
Model output & salt-and-pepper inputs & occlusion inputs \\
\midrule
clean image (reference) & 100\% & 100\% \\
corrupted input & 91\% & 77\% \\
universal AE (Task 1) & 97\% & 96\% \\
specialist, first run ($\alpha\approx0.3$) & 43\% & 79\% \\
specialist, final ($\alpha\approx0.6$) & 89\% & 88\% \\
soft MoE, first run ($\alpha\approx0.3$) & 46\% & 94\% \\
soft MoE, final ($\alpha\approx0.6$) & 97\% & 94\% \\
\bottomrule
\end{tabular}
}
\end{table}

\begin{table*}[t]
\centering
\caption{Hard routing on the test protocol: SSIM and L1 of oracle versus predicted routing, and the share of mis-routed inputs.}\label{tab:routing}
\small
\adjustbox{max width=\textwidth}{\begin{tabular}{lccccccc}
\toprule
 & \multicolumn{3}{c}{SSIM} & & \multicolumn{2}{c}{L1} & \\
\cmidrule(lr){2-5}\cmidrule(lr){6-7}
Input type & no restoration & universal & oracle & predicted & oracle & predicted & misroute \% \\
\midrule
clean & 1.000 & 0.814 & 1.000 & 0.996 & 0.0000 & 0.0010 & 2.34 \\
salt-and-pepper & 0.382 & 0.805 & 0.793 & 0.793 & 0.0444 & 0.0444 & 0.00 \\
Gaussian blur & 0.814 & 0.790 & 0.780 & 0.780 & 0.0449 & 0.0449 & 0.01 \\
occlusion & 0.713 & 0.707 & 0.689 & 0.690 & 0.0597 & 0.0594 & 1.19 \\
\midrule
all (36,690) & 0.672 & 0.772 & 0.779 & 0.779 & 0.0447 & 0.0447 & 0.59 \\
\bottomrule
\end{tabular}
}
\end{table*}

\textbf{Routing results.} Predicted routing (SSIM \nHardPredSSIM{}) is indistinguishable from oracle routing (\nHardOracleSSIM{}) because the classifier is accurate (Table~\ref{tab:routing}). Its advantage over the universal model (\nTOneSSIMall{}) comes entirely from clean images (\nHardSSIMclean{} against \nTOneSSIMclean{}). On corrupted inputs the specialists are slightly \emph{worse} than the universal model: \nHardSSIMsaltpepper{} against \nTOneSSIMsaltpepper{} for noise, \nHardSSIMblur{} against \nTOneSSIMblur{} for blur and \nHardSSIMocclusion{} against \nTOneSSIMocclusion{} for occlusion. A specialist has the same capacity and bottleneck as the universal model, so specialisation adds no reconstruction ability, and one shared search for three models and a batch of 32 (half the updates per epoch) may explain the small gap. The measured benefit of routing is the identity bypass.

\begin{figure*}[t]
\centering
\includegraphics[width=0.52\textwidth]{examples_hard.pdf}
\caption{Twelve test inputs for hard-routed restoration (predicted routing).}\label{fig:ex2}
\end{figure*}

\begin{figure}[t]
\centering
\includegraphics[width=0.85\columnwidth]{misroute_hard.pdf}
\caption{Classifier-induced failures: the four mis-routed clean images with the largest SSIM loss; predicted route next to oracle route.}\label{fig:mis}
\end{figure}

\begin{figure*}[t]
\centering
\includegraphics[width=0.7\textwidth]{curves_task2.pdf}
\caption{Task 2 training curves: classifier loss and accuracy, and validation PSNR of the three specialists in the first (dotted) and final (solid) run.}\label{fig:curves2}
\end{figure*}

\textbf{Classifier-caused failures.} The \nMisClean{} clean images sent to a blur or occlusion expert are real failures (SSIM \nMisCleanLoss{} on average relative to the oracle route). In Fig.~\ref{fig:mis} a clean dog on grass is taken for blurred and ``restored'' into a smeared brown scene (0.51 instead of 1.00), and an image with black letterbox bars and a cat on a black background are taken for occluded, so the occlusion expert invents content for the dark regions (0.54, 0.57). Conversely, 128 occlusions labelled ``clean'' (126 of them 10\% masks) are left untouched; this scores about $+0.12$ SSIM \emph{higher}, but only because SSIM ignores the black rectangle, which stays visible. Overall, mis-routing changes mean SSIM by only $+0.0001$ because the two groups cancel. The four worst outputs of the system are, as for Task~1, 35\% occlusions that hide the animal (SSIM 0.28--0.33); they are not routing errors, and the occlusion specialist fills the hole with a brownish patch even where the surroundings are green.
